\Needspace{75pt}
\pdfbookmark[1]{Berechnung der Mittelwerte der Beiträge und Pensionen}{witwen-mittelwerthe}
\begin{center}[\textit{Berechnung der Mittelwerte der Beiträge und Pensionen.}\end{center}
\noindent\hspace*{\parindent}Die Anzahl der Mitglieder, die in ununterbrochener Reihenfolge nach der Nummerierung ihrer Beitrittszeit bis zum 1. Oktober 1851 verstorben oder auf andere Weise aus der Professorenwitwenkasse ausgeschieden sind, beträgt 134. Für jedes dieser 134 ausgeschiedenen Mitglieder ist in der folgenden Tabelle zusammengestellt: das Jahr, an dessen 1. Okt. (und nur ausnahmsweise an dessen 1. Apr. bei Nr. 134) der Beitritt erfolgte, die Anzahl der Jahre und Monate, während derer die Beiträge oder die Witwen- und Waisenpensionen ein- oder ausgezahlt wurden, ebenso die bis zu dem Zeitpunkt verstrichenen Jahre und Monate, von dem an die Pension gerechnet wird. Aus diesen Daten sind die Werte in den danebenstehenden Spalten bestimmt worden. Sie sind zum Zeitpunkt des Eintritts bei einem Zinsfuß von \(3\tfrac{1}{2}\) Proz. und 4 Proz. den jährlich mit 10 Thl. eingezahlten Beiträgen und den mit 250 Thl. ausgezahlten Witwen- und Waisenpensionen gleichwertig. Die Ansätze von 10 Thl. und 250 Thl. hat \textsc{Gauss} wohl deshalb gewählt, weil der größte Teil der Tafel hierfür schon im Jahr 1845 berechnet war. Für die am 1. Okt. 1851 noch lebenden Witwen der Mitglieder Nr. 86. 103. 104. 109. 112. 114. 122 und 134 ist der Kapitalwert der nach diesem Zeitpunkt noch fälligen Pensionsbezüge mithilfe der Sterblichkeitstafeln von \textsc{Brune} bestimmt worden, unter der Voraussetzung, dass die Zahlung am Ende jedes Jahres, aber nur bis zum Schluss des Sterbemonats erfolgt. Die aus dieser Zusammenstellung unmittelbar hervorgehenden mittleren Kapitalwerte der Beiträge und Pensionen sind in der Bilanzrechnung von 1851 verwendet worden, nachdem durch Multiplikation mit einem Korrekturfaktor berücksichtigt worden ist, dass die Pensionen nicht jährlich, sondern halbjährlich gezahlt werden. Für die Waisenpensionen ist hier angenommen, dass sie bis zum Schluss des Monats galten, in dem das Kind sein 20\textsuperscript{ste} Lebensjahr vollendet hatte.

Die handschriftliche Tafel enthält, anders als der vorliegende Abdruck, keine Überschriften der einzelnen Spalten.]
